\section*{Limitations}
The comparison evaluates one recorded realization of each configuration on two related instruction-tuned models. It does not quantify seed variance, establish statistical superiority, or exhaust each method's hyperparameters. Budgets, inference-time property calls, and training requirements are not equalized. In particular, s1-style budget forcing omits s1K training, and one-revision Self-Refine does not evaluate longer refinement trajectories. Four-shot ICL is a contextual baseline. The benchmark therefore concerns these declared adaptations, not complete reproductions of every original method or causal comparisons of isolated mechanisms.

Success concerns frozen computational property evaluators. Using their success predicate as an inference-time verifier makes Sampling's terminal success equal its candidate-prefix coverage; this should not be interpreted as validation against independent experimental measurements. Strict directional improvement can be small, and the main metric imposes neither a minimum improvement magnitude nor a structural-similarity constraint. It does not establish synthesizability, biological efficacy, or clinical utility. The IND/OOD tasks follow property-combination splits whose source filters also differ, limiting conclusions about generalization. Generation workload omits unreturned beam-search and property-tool computation. Shared difficulty groups are defined by Mistral reference runs, not intrinsic difficulty; conditional quality metrics average different successful subsets across methods.

\section*{Ethical Considerations}
This study evaluates computational molecular optimization using existing datasets and models. Property predictions and generated structures are not experimentally validated, and optimization of the selected proxies can overlook unmeasured toxicity or other undesirable effects. The benchmark should be used to study inference procedures, with domain-specific validation required before downstream scientific or biomedical use. No clinical conclusions or recommendations follow from the reported success rates.
